\documentclass[11pt]{article}
\usepackage[letterpaper,margin=1in]{geometry}
\usepackage[T1]{fontenc}
\usepackage{lmodern,microtype,amsmath,amssymb,amsthm}
\usepackage{xcolor}
\usepackage[colorlinks=true,linkcolor=black,citecolor=blue!45!black,urlcolor=blue!45!black]{hyperref}
\usepackage{xurl}
\newtheorem{theorem}{Theorem}[section]
\newtheorem{lemma}[theorem]{Lemma}
\newtheorem{corollary}[theorem]{Corollary}
\theoremstyle{definition}
\newtheorem{definition}[theorem]{Definition}
\theoremstyle{remark}
\newtheorem{remark}[theorem]{Remark}
\newcommand{\ED}{\mathsf{ED}}
\newcommand{\vol}{\mathrm{vol}}
\newcommand{\F}{\mathbb{F}}
\newcommand{\SIZE}{\mathsf{SIZE}}
\newcommand{\NTIME}{\mathsf{NTIME}}
\newcommand{\coNTIME}{\mathsf{coNTIME}}
\newcommand{\Lclass}{\mathsf{L}}
\newcommand{\ENP}{\mathsf{E}^{\mathsf{NP}}}
\hypersetup{pdftitle={Element distinctness requires superlinear circuits, assuming a tensor-network volume bound},
  pdfauthor={Charlie Cruz},
  pdfkeywords={circuit lower bounds, element distinctness, tensor networks, volume, superlinear lower bounds, Boolean circuits}}
\title{Element Distinctness Requires Superlinear Circuits,\texorpdfstring{\\}{ }Assuming a Tensor-Network Volume Bound}
\author{Charlie Cruz}
\date{November 24, 2026}
\begin{document}
\maketitle
\begin{abstract}
\emph{Element distinctness} asks whether a list of words contains a
repeated word. A \emph{tensor network} represents a function as a sum,
over the values of its internal wires, of products of small tables, and
its \emph{volume} counts its vertices and wires. The recent paper
\emph{The matrix multiplication exponent is nine fourths} \cite{OAI}
states a superlinear lower bound on the volume of tensor networks for
list distinctness (its Theorem~2.2). We observe that this theorem implies
a superlinear lower bound for element distinctness against unrestricted
Boolean circuits. Assuming it, there are constants $\sigma,c>0$ such
that, for all sufficiently large $n$, every fan-in-two circuit over the
full binary basis computing element distinctness on $n$ bits has at
least $c\,n(\log n)^{\sigma}$ gates. Element distinctness is in logspace
and its complement is in nondeterministic linear time, so
$\Lclass\cap\coNTIME[n]\not\subseteq\SIZE[O(n(\log n)^{\delta})]$ for all
$\delta<\sigma$. The deduction is short: a circuit converts into a tensor
network of linear volume. All of the difficulty lies in the imported
theorem, and we state which parts of its proof we have checked.
\end{abstract}

\section{Introduction}

\subsection{Background}
Proving that an explicit Boolean function requires superlinear circuit
size is a basic open problem in complexity theory. It has not been ruled
out that every function in $\mathsf{NP}$ has circuits of size $O(n)$. The
best lower bound for an explicit function against fan-in-two circuits
over the full binary basis has improved from $3n-o(n)$ \cite{Blum} to
$(3+\tfrac{1}{86})n$ \cite{FGHK} and then to about $3.1n$ \cite{LY}.
These bounds come from gate elimination, and they are far from even
$n\log\log n$.

\subsection{Results}
The November 24, 2026 release on matrix multiplication \cite{OAI}
contains a statement strong enough to give a superlinear bound. The
statement we use is Theorem~2.2 of \cite{OAI}, an intermediate theorem
about tensor networks that test whether a list of words has a repeat; we
do not use the main result $\omega=9/4$. That theorem is stated without
hypotheses. We restate the case we need as Theorem~\ref{thm:import} in
Section~\ref{sec:prelim}, where we also define $\ED_n$, element
distinctness on $n$ input bits (Definition~\ref{def:ed}).

\begin{theorem}\label{thm:main}
Assume Theorem~\ref{thm:import} (Theorem~2.2 of \cite{OAI}). There are
constants $\sigma,c>0$ such that for every sufficiently large $n$, every
fan-in-two Boolean circuit computing element distinctness $\ED_n$ on $n$
input bits has more than $c\,n(\log n)^{\sigma}$ gates.
\end{theorem}

Element distinctness is an easy problem. A logspace machine decides it
by comparing all pairs of words with a few counters, and a
nondeterministic machine decides its complement by guessing a repeated
pair and checking it (Lemma~\ref{lem:uniform}).

\begin{corollary}\label{cor:classes}
Assume Theorem~\ref{thm:import}. For every fixed $0\le\delta<\sigma$,
\[
\Lclass\cap\coNTIME[n]\not\subseteq\SIZE[O(n(\log n)^{\delta})].
\]
Moreover, the lower bound holds at every sufficiently large input length,
not just infinitely often.
\end{corollary}

For comparison, the refutation of ETH \cite{ETH} gives an
$\Omega(n\sqrt{\log n})$ lower bound for a language in $\ENP$, but only at
infinitely many input lengths \cite[Corollary~1.4]{ETH}. Its exponent of
$\log n$ is larger than ours if $\sigma<1/2$, but its hard function lies
in a much larger class. To our knowledge, Theorem~\ref{thm:main} is the
first superlinear lower bound against unrestricted circuits for any
function in $\mathsf{P}$.

\subsection{Techniques}
The proof uses one idea, which already appears in \cite{OAI}: a Boolean
circuit is a tensor network. Put a $0/1$ table at every gate that checks
the gate's input-output relation, put an equality table wherever a wire
fans out, and pin the output to $1$. Summing over all values of the
internal wires then leaves exactly one nonzero term, namely the
evaluation of the circuit. So the network represents the circuit's
function exactly, over any field, and its volume is linear in the size
of the circuit (Lemma~\ref{lem:convert}). In the proof of
\cite[Proposition~2.3]{OAI}, this conversion is applied to a sorting
network to bound the volume of list distinctness from above. We apply
the same conversion in the other direction, against the lower bound of
\cite[Theorem~2.2]{OAI}.

\subsection{Status of the imported theorem}
Theorem~\ref{thm:main} is exactly as strong as Theorem~\ref{thm:import},
and the proof of Theorem~\ref{thm:import} is long. It goes through a case
analysis on whether the field has \emph{ports}, a nonlinear Poincar\'e
property of certain metric spaces attached to tensors \cite{Geom}. The
two cases use different tools: the CAT(0) geometry of spectral points in
one case, and a self-referential counting argument from \cite{HC} in the
other. We have read the parts of the argument that interface with list
distinctness and found no gap. Section~\ref{sec:status} lays out what has
and has not been checked. We think the reader should treat
Theorem~\ref{thm:main} as one treats any striking corollary of a new and
long proof: it is a reason to scrutinize that proof, and a reason to care
about the answer.

The conversion appears one paragraph after Theorem~2.2 in \cite{OAI}, but
\cite{OAI} does not state this consequence. We do not know why.

\section{Preliminaries}\label{sec:prelim}
Circuits are single-output, with fan-in at most two, arbitrary binary
Boolean gates, and unrestricted fan-out. Size counts gates. All
logarithms are base two. Time bounds refer to multitape Turing machines.

\begin{definition}\label{def:ed}
For $w\ge1$ let $n_w=2w\cdot2^w$, and let $D_w:\{0,1\}^{n_w}\to\{0,1\}$
take a list of $2^w$ words of $2w$ bits each and return $1$ exactly when
the words are pairwise distinct. For every $n\ge n_1$ let $w(n)$ be the
largest $w$ with $n_w\le n$. Define $\ED_n:\{0,1\}^n\to\{0,1\}$ to apply
$D_{w(n)}$ to the first $n_{w(n)}$ input bits and ignore the rest, and
let $\ED=\bigcup_n\ED_n^{-1}(1)$.
\end{definition}

Any reasonable encoding of element distinctness works equally well. All
we need is that $D_w$ appears inside it, after fixing some input bits to
constants, at input length $O(n_w)$.

\paragraph{Tensor networks.}
We follow \cite[Section~2.1]{OAI}. Fix a finite field $k$. A \emph{bit
tensor} of arity $u$ is a function $\{0,1\}^u\to k$. A tensor network
places a bit tensor (a \emph{table}) at each vertex. Each wire pairs two
argument positions, and each unpaired position is a named exterior bit.
For a setting of the exterior bits, multiply the vertex entries and sum
over all settings of the internal wires. The result is the tensor the
network represents. The network is \emph{trivalent} if every vertex has
arity at most three. Its \emph{volume} counts vertices and all wires,
including exterior wires. Let $\vol_k(T)$ be the minimum volume of a
trivalent network over $k$ representing $T$. We view $D_w$ as a bit
tensor with values $0,1\in k$.

\begin{theorem}[{\cite[Theorem~2.2]{OAI}}, case $t=1$]\label{thm:import}
For every finite field $k$ there is $\sigma>0$ such that for all
sufficiently large $w$, $\vol_k(D_w)>n_w w^{\sigma}$.
\end{theorem}

The theorem in \cite{OAI} is stronger. It handles $t$ independent lists
at once, uniformly for $1\le t\le 2^w$, even when the network shares work
across lists. We need only one list.

\section{Upper bounds for element distinctness}
\begin{lemma}\label{lem:uniform}
$\ED\in\Lclass$, and the complement of $\ED$ is in $\NTIME[n]$.
\end{lemma}
\begin{proof}
From $n$ one computes $w=w(n)$ and $n_w$ by arithmetic on
$O(\log n)$-bit numbers.

For logspace, loop over all pairs $i<j$ of word indices with two
$O(w)$-bit counters, and compare the two words bit by bit with a third
counter. Accept exactly when no pair agrees.

For the complement, the only care needed is to stay within linear time
on tapes. Obtain $n$ in binary with an amortized constant-time counter,
and compute $w$ and $n_w$ in polylogarithmic time. Write a marked block
of length $2w$ on a work tape. Moving a head cyclically along this block
while scanning the input locates word boundaries at constant cost per
input bit. A second counter, decremented with amortized constant cost,
stops the scan at position $n_w$. During one scan, guess the start of a
word $i$ at a boundary and copy its $2w$ bits to a work tape. Then guess
the start of a later word $j$ and compare it with the copy. Accept if
they agree.
\end{proof}

For circuits, a sorting network followed by comparisons of adjacent
words gives circuits for $D_w$ of size $O(n_w w^2)$ \cite{OAI}. So
Theorem~\ref{thm:main} is tight up to the exponent of $\log n$.

\section{Circuits as tensor networks}
The following conversion is from \cite[proof of Proposition~2.3]{OAI}; we
write out the size accounting.

\begin{lemma}\label{lem:convert}
Let $C$ be a circuit with $s$ gates on $n$ inputs computing
$f:\{0,1\}^n\to\{0,1\}$. Over every field $k$ there is a trivalent network
with exterior bits $x_1,\dots,x_n$ representing $f$, with values $0,1\in k$,
and with volume at most $16s+4n+8$.
\end{lemma}
\begin{proof}
Call the inputs and gate outputs \emph{signals}. Give each signal a set
of \emph{uses}: one for each gate input that reads it, plus one extra use
for the output signal of $C$. The total number of uses is at most
$2s+1$.

Put the following tables in the network.
\begin{itemize}
\item For each gate $g$ with inputs $a,b$ and output $c$, the table
$T_g(a,b,c)=[c=g(a,b)]$. Gates of fan-in one or zero get tables of
arity two or one in the same way.
\item For the output signal, a pin $P(c)=[c=1]$ attached as its extra use.
\item For each signal, a distribution structure joining its source (the
exterior bit for an input, the output position of $T_g$ for a gate) to
its uses. Let $d$ be the number of uses. If $d=1$, use a single wire.
If $d\ge2$, use a path of $d-1$ equality tables
$E(x,y,z)=[x=y=z]$. If $d=0$, end the source at a table that is
identically $1$.
\end{itemize}
Fix an exterior setting $x$. The equality tables force all copies of a
signal to agree, and the gate tables force each gate output to equal its
gate's value on its inputs. So exactly one internal assignment has a
nonzero product: the evaluation of $C$ on $x$. The ones-tables and the
pin then contribute the factor $[C(x)=1]$. Hence the network represents
$f$, with entries $0$ and $1$ in $k$, in every characteristic.

There are $s$ gate tables, one pin, at most $n+s$ ones-tables, and at most
$2s+1$ equality tables in total. That gives at most $V\le 4s+n+2$
vertices. Every wire has at least one endpoint at a vertex, and every
vertex has arity at most three, so there are at most $3V$ wires. The
volume is at most $4V\le 16s+4n+8$.
\end{proof}

\section{Proofs of the main results}
\begin{proof}[Proof of Theorem~\ref{thm:main}]
Let $n$ be large, put $w=w(n)$, and let $C$ be a circuit for $\ED_n$ with
$s$ gates. Apply Lemma~\ref{lem:convert} to $C$ over $k=\F_2$. Then attach
a pin $Z(x)=[x=0]$ to each of the last $n-n_w$ exterior bits. This turns
those exterior wires into internal wires and adds $n-n_w$ vertices. We
pin inside the network, rather than fixing inputs of $C$, to avoid
handling gates with constant inputs. Since $\ED_n$ ignores these bits,
the pinned network represents $D_w$ on the remaining $n_w$ exterior bits.

By maximality of $w$ we have $n<n_{w+1}=2n_w(w+1)/w\le4n_w$. Comparing
with Theorem~\ref{thm:import},
\[
16s+5n+8\;\ge\;\vol_{\F_2}(D_w)\;>\;n_w w^{\sigma},
\qquad\text{so}\qquad
s>\frac{n_w w^{\sigma}-20n_w-8}{16}\ge\frac{n_w w^{\sigma}}{32}
\]
for large $w$. Also $2(w+1)\le2^{w+1}$, so $n<n_{w+1}\le 2^{2(w+1)}$.
This gives $w>\tfrac12\log n-1\ge\tfrac14\log n$ for large $n$. Therefore
\[
s>\frac{n_w w^{\sigma}}{32}>\frac{n}{128}\Bigl(\frac{\log n}{4}\Bigr)^{\sigma},
\]
and we may take $c=4^{-\sigma}/128$.
\end{proof}

\begin{proof}[Proof of Corollary~\ref{cor:classes}]
Lemma~\ref{lem:uniform} places $\ED$ in $\Lclass\cap\coNTIME[n]$. For
$\delta<\sigma$ and any constant $C'$, we have
$C'n(\log n)^{\delta}<c\,n(\log n)^{\sigma}$ for large $n$. So no circuit
family of size $O(n(\log n)^\delta)$ computes $\ED$, even at infinitely
many lengths.
\end{proof}

\section{Status of Theorem~\ref{thm:import}}\label{sec:status}
Everything above is elementary. The conversion is already written down
in \cite{OAI}, and Theorem~\ref{thm:import} is stated there with no
hypothesis. In particular it assumes neither $\omega<9/4$ (unlike the
compression theorem it is later combined with) nor Weak ETH. So the
question is entirely whether Theorem~\ref{thm:import} is correct. We
outline its proof and state what we know about each piece.

\paragraph{From distinctness to routing.}
Two calls to a network for $D_w$ can check a \emph{tagged permutation}:
that a list of tags forms a permutation $\pi$ and that each payload has
been delivered to position $\pi(i)$ \cite[Section~4.1]{OAI}. Now pin the
tags to a uniformly random permutation. What remains is a fixed network
topology that must route $r$ payload bits, each to a uniformly random
destination among $q$ positions. We checked this reduction, including
over $\F_2$.

\paragraph{Case 1: no ports.}
\cite[Lemma~4.1]{OAI} shows that any such topology has size
$\Omega(r\sqrt{\log q})$. The argument places the network in a spectral
space whose distances behave like wire lengths. It then compares the
total distance the payloads must travel with the total length of the
wires, using a graph entropy bound \cite[Corollary~5.2]{HC} and a minimax
step. Plain counting gives nothing superlinear here, because the pinned
tags carry exactly enough information to specify the permutation.
Everything above linear comes from the geometry. We checked
\cite[Lemma~4.1]{OAI}. However, it rests on results about CAT(0)
spectral spaces that we have only spot-checked, namely
\cite[Theorem~4.3]{Geom} and \cite[Theorems~4.2, 5.1 and~6.1]{HC}. A
concrete test would be a spectral point whose wire distances on four
sites violate the CAT(0) quadrilateral inequality. Such a point would
break \cite[Lemma~4.1]{OAI}.

\paragraph{Case 2: ports.}
Here a cheap network for $D_w$ is used, through the descent of
\cite[Proposition~14.1]{Geom}, to compress batches of a signed counting
tensor $U^*_\ell$ below the exact batch lower bound of
\cite[Theorem~12.1]{HC}. That theorem is stated without hypothesis, but
its proof (Sections~13--15 of \cite{HC}) assumes ports. This is
harmless, since it is only invoked in the ports case. Its proof is a
diagonal argument, in which a program uses the compression to evaluate
its own success count and add one. We checked every parameter step of
the ports case in \cite{OAI}. The diagonal argument in \cite{HC} deserves
a careful line-by-line audit of its clocks and budgets.

\paragraph{Which case holds.}
We do not know. The multiplication paper shows that ports in
characteristic zero would refute Weak ETH \cite[Theorem~7.1]{Mult}. That
is suggestive, but our field is $\F_2$, and that theorem says nothing
about characteristic two. So both cases need to be right.

\paragraph{Comparison with network coding.}
\cite[Lemma~4.1]{OAI} is, in effect, an unconditional superlinear lower
bound for bounded-degree networks that route bits to data-dependent
destinations, with arbitrary coding allowed along the way. For Boolean
circuits, bounds of this kind were previously known only under the
network coding conjecture \cite{AFKL}. That is a good reason to look hard
at \cite[Lemma~4.1]{OAI}. By itself it is no reason to doubt the lemma,
since the spectral machinery of \cite{Geom} is a new tool for charging
coded information for the distance it travels.

\section{Discussion}
\paragraph{Open problems.}
\begin{itemize}
\item \emph{Make $\sigma$ explicit.} If $\F_2$ has no ports, the proof
gives every $\sigma<1/2$. What is the best exponent one can extract from
the ports case? Is $\Omega(n\log n)$ possible for element distinctness,
matching the sorting upper bound up to a $\log n$ factor?
\item \emph{Polynomially superlinear bounds.} Can the volume machinery
of \cite{OAI,Geom,HC} give an $n^{1+\varepsilon}$ lower bound for any
explicit function? Hardness magnification \cite{MMW,OS} shows that
$n^{1+\varepsilon}$ lower bounds for certain compression problems would
already separate $\mathsf{NP}$ from $\mathsf{P}/\mathsf{poly}$. Element
distinctness is not one of those problems, so no magnification applies
directly to Theorem~\ref{thm:main}. It would be interesting to know
whether the tensor-network method can reach them.
\item \emph{Other functions.} The argument needs only that the function
forces data-dependent routing. Sorting (as a multi-output function),
checking a permutation, and set disjointness on structured inputs are
natural candidates for the same treatment.
\item \emph{Which case holds for $\F_2$?} Deciding whether characteristic
two admits ports would remove half of the proof from consideration.
\end{itemize}

\paragraph{Acknowledgments.}
The bound is an immediate corollary of \cite[Theorem~2.2]{OAI}, and all
credit for the underlying theorem belongs to its authors.

\begin{thebibliography}{99}
\bibitem{AFKL}
P.~Afshani, C.~B. Freksen, L.~Kamma, and K.~G. Larsen. Lower bounds for
multiplication via network coding. In \emph{Proceedings of ICALP 2019},
LIPIcs 132, pages 10:1--10:12, 2019.

\bibitem{Blum}
N.~Blum. A Boolean function requiring $3n$ network size.
\emph{Theoretical Computer Science}, 28(3):337--345, 1984.

\bibitem{FGHK}
M.~G. Find, A.~Golovnev, E.~A. Hirsch, and A.~S. Kulikov. A better-than-$3n$
lower bound for the circuit complexity of an explicit function.
In \emph{Proceedings of FOCS 2016}, pages 89--98, 2016.

\bibitem{LY}
J.~Li and T.~Yang. $3.1n-o(n)$ circuit lower bounds for explicit functions.
In \emph{Proceedings of STOC 2022}, pages 1180--1193, 2022.

\bibitem{MMW}
D.~M. McKay, C.~D. Murray, and R.~R. Williams. Weak lower bounds on
resource-bounded compression imply strong separations of complexity
classes. In \emph{Proceedings of STOC 2019}, pages 1215--1225, 2019.

\bibitem{OS}
I.~C. Oliveira and R.~Santhanam. Hardness magnification for natural
problems. In \emph{Proceedings of FOCS 2018}, pages 65--76, 2018.

\bibitem{OAI}
OpenAI. The matrix multiplication exponent is nine fourths.
Preprint, November 24, 2026.
\href{https://github.com/openai/math/blob/3344da59991ed94f9ba711217b29c1d760badfa6/preprints/The-matrix-multiplication-exponent-is-nine-fourths-November-24-2026/paper.pdf}{Manuscript}.

\bibitem{HC}
OpenAI. Circuit lower bounds for Hamiltonian cycle and the permanent.
Preprint, November 24, 2026.
\href{https://github.com/openai/math/blob/3344da59991ed94f9ba711217b29c1d760badfa6/preprints/Circuit-lower-bounds-for-Hamiltonian-cycle-and-the-permanent-November-24-2026/paper.pdf}{Manuscript}.

\bibitem{Geom}
OpenAI. Euclidean geometry in the asymptotic spectrum of tensors.
Preprint, November 24, 2026.
\href{https://github.com/openai/math/blob/3344da59991ed94f9ba711217b29c1d760badfa6/preprints/Euclidean-geometry-in-the-asymptotic-spectrum-of-tensors-November-24-2026/paper.pdf}{Manuscript}.

\bibitem{Mult}
OpenAI. The logarithmic exponent of integer multiplication on tapes.
Preprint, November 24, 2026.
\href{https://github.com/openai/math/blob/3344da59991ed94f9ba711217b29c1d760badfa6/preprints/The-logarithmic-exponent-of-integer-multiplication-on-tapes-November-24-2026/paper.pdf}{Manuscript}.

\bibitem{ETH}
OpenAI. The Exponential Time Hypothesis is False.
Preprint, October 25, 2026.
\href{https://github.com/openai/math/tree/3344da59991ed94f9ba711217b29c1d760badfa6/preprints/The-Exponential-Time-Hypothesis-is-False-October-25-2026}{Manuscript repository}.
\end{thebibliography}
\end{document}
